%
% file: dedication.tex
% author: V?ctor Bre?a-Medina
% description: Contains the text for thesis dedication
%


% ------------------------------------------------------------------
% Acronym setup
%   * long-short style: the first use of an acronym in the document
%     prints "long form (SHORT)", every later use prints just "SHORT".
%   * \newacro works exactly like \newacronym but ALSO defines a
%     \SHORT command for the acronym, e.g. \CMB -> \gls{CMB_}.
%     \xspace keeps the trailing space ("\CMB photons" -> "CMB photons")
%     while still suppressing it before punctuation ("\CMB." -> "CMB.").
%     For capitalised (sentence start) or plural forms use the
%     glossaries commands directly: \Gls{key}, \glspl{key}, \Glspl{key}.
% ------------------------------------------------------------------
\RequirePackage{xspace}
\setacronymstyle{long-short}
\newcommand{\newacro}[4][]{%
  \newacronym[#1]{#2}{#3}{#4}%
  \expandafter\newcommand\csname #3\endcsname{\gls{#2}\xspace}%
}

% define many of the acronyms (all via \newacro, so each gets a \SHORT command)
% Long forms are lowercase except genuine proper nouns (instrument names,
% person-names like Compton/Newtonian, the Big Bang, the TensorFlow brand).
\newacro{DM_}{DM}{dark matter}
\newacro{DGRB_}{DGRB}{diffuse gamma-ray background}
\newacro{EGB_}{EGB}{extragalactic gamma-ray background}
\newacro{LAT_}{LAT}{Large Area Telescope}
\newacro{SCD_}{SCD}{source-count distribution}
\newacro[longplural={active galactic nuclei}]{AGN_}{AGN}{active galactic nucleus}
\newacro{UV_}{UV}{ultraviolet}
\newacro{IC_}{IC}{inverse Compton}
\newacro[longplural={flat-spectrum radio quasars}]{FSRQ_}{FSRQ}{flat-spectrum radio quasar}
\newacro[longplural={misaligned active galactic nuclei}]{mAGN_}{mAGN}{misaligned active galactic nucleus}
\newacro[longplural={star-forming galaxies}]{SFG_}{SFG}{star-forming galaxy}
\newacro[longplural={cosmic rays}]{CR_}{CR}{cosmic ray}
\newacro{NN_}{NN}{neural network}
\newacro{ML_}{ML}{machine learning}
\newacro{LSP_}{LSP}{lightest supersymmetric particle}
\newacro[longplural={weakly interacting massive particles}]{WIMP_}{WIMP}{weakly interacting massive particle}
\newacro{PSF_}{PSF}{point-spread function}
\newacro{IRF_}{IRF}{instrument response function}
\newacro{TPU_}{TPU}{tensor processing unit}
\newacro{MLNN_}{MLNN}{multilayer neural network}
\newacro{CNN_}{CNN}{convolutional neural network}
\newacro{TF_}{TF}{TensorFlow}

\newacro{resnet_}{ResNet}{residual neural network}
\newacro{ICS_}{ICS}{inverse Compton scattering}
\newacro[longplural={supernova remnants}]{SNR_}{SNR}{supernova remnant}

% Chapter 1 (cosmology / evidence for dark matter)
\newacro{CMB_}{CMB}{cosmic microwave background}
\newacro{BBN_}{BBN}{Big Bang nucleosynthesis}
\newacro{SZ_}{SZ}{Sunyaev-Zel'dovich}
\newacro{ICM_}{ICM}{intracluster medium}
\newacro[longplural={dwarf spheroidal galaxies}]{dSph_}{dSph}{dwarf spheroidal galaxy}
\newacro{BAO_}{BAO}{baryon acoustic oscillations}
\newacro{MOND_}{MOND}{modified Newtonian dynamics}
\newacro[longplural={primordial black holes}]{PBH_}{PBH}{primordial black hole}
\newacro[longplural={feebly interacting massive particles}]{FIMP_}{FIMP}{feebly interacting massive particle}
\newacro{WDM_}{WDM}{warm dark matter}

%\newacronym{DM}{Dark Matter}
\newacro{UGRB_}{UGRB}{unresolved gamma-ray background}


\newacro{CTAO_}{CTAO}{Cherenkov Telescope Array Observatory}
\newacro{SBI_}{SBI}{simulation-based inference}

% Chapter 1 (particle candidates / indirect detection)
\newacro{LHC_}{LHC}{Large Hadron Collider}
\newacro{ALP_}{ALP}{axion-like particle}
\newacro{VIB_}{VIB}{virtual internal bremsstrahlung}
\newacro{SIDM_}{SIDM}{self-interacting dark matter}
\newacro{ISRF_}{ISRF}{interstellar radiation field}
\newacro{NFW_}{NFW}{Navarro-Frenk-White}

% Chapter 2 (gamma-ray sky / Fermi-LAT)
\newacro{ISM_}{ISM}{interstellar medium}
\newacro{GDE_}{GDE}{Galactic diffuse emission}
\newacro{MSP_}{MSP}{millisecond pulsar}
\newacro{ACD_}{ACD}{anticoincidence detector}
\newacro{SED_}{SED}{spectral energy distribution}
\newacro{EDISP_}{EDISP}{energy dispersion}

% Chapter 3 (statistical methods / simulation-based inference)
\newacro{MLE_}{MLE}{maximum likelihood estimator}
\newacro{TS_}{TS}{test statistic}
\newacro{KS_}{KS}{Kolmogorov--Smirnov}
\newacro{KL_}{KL}{Kullback--Leibler}
\newacro{AIC_}{AIC}{Akaike information criterion}
\newacro{MCMC_}{MCMC}{Markov chain Monte Carlo}
\newacro{NPE_}{NPE}{neural posterior estimation}
\newacro{NRE_}{NRE}{neural ratio estimation}
\newacro{NLE_}{NLE}{neural likelihood estimation}
\newacro{ABC_}{ABC}{approximate Bayesian computation}
\newacro{SBC_}{SBC}{simulation-based calibration}
\newacro{TARP_}{TARP}{tests of accuracy with random points}
\newacro{SGD_}{SGD}{stochastic gradient descent}
\newacro{KDE_}{KDE}{kernel density estimation}
\newacro{APS_}{APS}{angular power spectrum}

% Chapter 4 (Galactic Center excess)
\newacro{GCE_}{GCE}{Galactic Center excess}
\newacro{NPTF_}{NPTF}{non-Poissonian template fitting}
\newacro[longplural={low-mass X-ray binaries}]{LMXB_}{LMXB}{low-mass X-ray binary}

% Chapter 7 (probabilistic cataloging)
\newacro{QF_}{QF}{quality factor}
\newacro[longplural={regions of interest}]{RoI_}{RoI}{region of interest}

% Chapter 8 (cross-correlations / CTAO)
\newacro{EBL_}{EBL}{extragalactic background light}
\newacro{IACT_}{IACT}{imaging atmospheric Cherenkov telescope}
\newacro{LST_}{LST}{Large-Sized Telescope}
\newacro{MST_}{MST}{Medium-Sized Telescope}
\newacro{EGAL_}{EGAL}{Extragalactic Survey}
\newacro{HSP_}{HSP}{high-synchrotron-peaked}
\newacro{LISP_}{LISP}{low-to-intermediate-synchrotron-peaked}

% gls-only entries: SHORT starts with a digit, so the auto-generated \SHORT
% command is unreachable; use \gls{KEY_}
\newacro{1pPDF_}{1pPDF}{1-point probability distribution function}
\newacro{2MASS_}{2MASS}{Two Micron All-Sky Survey}
\newacro{2MRS_}{2MRS}{2MASS Redshift Survey}

% signal-to-noise ratio shares the visible SHORT "SNR" with SNR_ (supernova remnant);
% raw \newacronym because a second \newacro would redefine \SNR. Use \gls{SNRatio_}.
\newacronym{SNRatio_}{SNR}{signal-to-noise ratio}




